\section{Related work}
\label{related}

Vidros et al.~\cite{vidros2017} introduced EMSCAD and classified it with bag-of-words features and
standard classifiers, on a balanced subset of the data. Most later studies combine text features
with meta-features such as a missing company profile and use linear, tree or ensemble classifiers;
Taneja et al.~\cite{taneja2025} list ten of them with 91\% to 99.9\% accuracy, including one measured
on a resampled test set~\cite{amaar2022}. Transformers have since been applied to the same data:
Akram et al.~\cite{akram2024} fine-tune BERT~\cite{devlin2019} and RoBERTa with SMOTE oversampling,
and Fraud-BERT~\cite{taneja2025} reports F1 0.93 against 0.78 for logistic regression, on one 80:20
split without deduplication or seeds.

Two further results shape our protocol. Oversampling such as SMOTE~\cite{chawla2002} is
legitimate during training but makes test numbers meaningless when applied to the test set, and
precision-recall curves are more informative than ROC curves when negatives
dominate~\cite{saito2015}. For explaining a neural model, Integrated
Gradients~\cite{sundararajan2017} attributes a prediction to the input tokens, with the guarantee
that the attributions add up to the change in output from a neutral reference input.

Taken together, EMSCAD results are hard to compare, rarely report their spread over splits, and
rarely isolate what a fine-tuned transformer adds over a tuned word-count baseline. This project
fixes one protocol and puts every model through it.
